% 04_complejidad.tex — Slide 4: Complejidad y NP-completitud
\begin{frame}{Complejidad: por que 3DM es NP-completo}

\begin{columns}[T]
\begin{column}{0.50\textwidth}
{\small
\textbf{Pertenencia a NP:}
\begin{itemize}
  \item Certificado: $k$ indices de tripletas
  \item Verificador: 3 bitmasks en $O(n)$
\end{itemize}

\smallskip
\textbf{NP-dificultad (Karp 1972):}\\
Reduccion $\text{3-SAT} \le_p \mathbf{3DM}$ via gadgets:
\begin{itemize}
  \item \textbf{G-V}: codifica variables
  \item \textbf{G-C}: codifica clausulas
  \item \textbf{G-G}: recolecta basura
\end{itemize}
Instancia: $n{=}6q$ elems., $m{=}O(q^3)$ trips.

\smallskip
\textbf{max-3DM:} APX-dificil
\begin{itemize}
  \item Sin PTAS salvo P=NP
  \item Mejor aprox.: $\approx 0.66$
\end{itemize}
}
\end{column}

\begin{column}{0.47\textwidth}

\begin{block}{Cadena de reducciones}
\centering
\smallskip
\begin{tikzpicture}[
  every node/.style={font=\small},
  arrow/.style={->, thick, mainblue}
]
  \node[draw, rounded corners, fill=blue!10, minimum width=2.8cm, minimum height=0.6cm]
        (sat)  at (0, 0)   {3-SAT (\#11)};
  \node[draw, rounded corners, fill=red!10,  minimum width=2.8cm, minimum height=0.6cm,
        line width=1.5pt]
        (tdm)  at (0,-1.3) {\textbf{3DM (\#17)}};
  \node[draw, rounded corners, fill=gray!10, minimum width=2.8cm, minimum height=0.6cm]
        (sp)   at (0,-2.6) {Set Packing (\#4)};
  \draw[arrow] (sat) -- (tdm) node[midway, right, font=\tiny] {Karp '72};
  \draw[arrow] (tdm) -- (sp)  node[midway, right, font=\tiny] {$\le_p$};
\end{tikzpicture}

\smallskip
{\small\textit{La 3$^{\mathrm{a}}$ dimension rompe la unimodularidad que hace P al matching bipartito.}}
\end{block}

\end{column}
\end{columns}

\end{frame}
